\section{Implementation Surface}
\label{sec:otel-weaver-exp-005-collector-to-intake-implementation}

\subsection{Source Files}
\label{sec:otel-weaver-exp-005-collector-to-intake-implementation-sources}

The EXP-005 implementation surface consists of three manufactured artifacts and one pending
output artifact:

\begin{description}
  \item[\texttt{harness.py}] The Python synthetic load generator. Fires three span classes
        using the OpenTelemetry Python SDK: Span~A (fully compliant, carries both
        \texttt{process.pi.instance\_id} and \texttt{process.pi.witness.hash}), Span~B
        (missing \texttt{process.pi.instance\_id}, expected to be dropped at \texttt{filter/pi}),
        and Span~C (carries \texttt{process.pi.instance\_id} but missing
        \texttt{process.pi.witness.hash}, expected to pass the collector but fail the
        downstream live-check gate). The harness also emits the
        \texttt{pi.feedstock.operation\_count} cumulative sum metric tagged with the instance
        ID, providing a counter that can be used to verify round-trip span delivery counts.
        Source: \texttt{experiments/exp-005-collector-to-pi-intake/harness.py}.

  \item[\texttt{config/otel-collector-config.yaml}] The authoritative collector configuration.
        Declares the full three-stage pipeline: OTLP receiver (gRPC \texttt{:4317}, HTTP
        \texttt{:4318}), \texttt{filter/pi} and \texttt{transform/pi} OTTL processors, and
        \texttt{file/pi\_json} exporter. Contains production-ready OTTL transform rules
        including the \texttt{set(attributes["pi.intake.ingested\_at"], ...)} and
        \texttt{set(attributes["pi.intake.collector\_version"], ...)} transform statements.
        Source: \texttt{experiments/exp-005-collector-to-pi-intake/config/otel-collector-config.yaml}.

  \item[\texttt{README.md}] The experiment specification. Contains the full pipeline diagram,
        the OCEL~2.0 JSON event schema, the Docker launch command, and the admissibility rules.
        This is the canonical human-readable specification for the experiment.
        Source: \texttt{experiments/exp-005-collector-to-pi-intake/README.md}.

  \item[\texttt{output/pi\_intake\_events.json}] The intended OCEL~2.0 JSON intake artifact.
        PARTIAL: the \texttt{output/} directory exists but \texttt{pi\_intake\_events.json}
        has not been materialised because no live Docker collector run has been executed.
\end{description}

\subsection{Commands}
\label{sec:otel-weaver-exp-005-collector-to-intake-implementation-commands}

The README specifies the following Docker launch command to start the collector:
\begin{verbatim}
docker run --rm \
  -v $(pwd)/config:/etc/otelcol-contrib \
  -v $(pwd)/output:/output \
  -p 4317:4317 -p 4318:4318 \
  otel/opentelemetry-collector-contrib:0.98.0 \
  --config /etc/otelcol-contrib/otel-collector-config.yaml
\end{verbatim}

Once the collector is running, the harness is executed with:
\begin{verbatim}
python harness.py
\end{verbatim}

The harness sends spans to the OTLP gRPC endpoint at \texttt{localhost:4317} by default.
These commands have not been executed against a live collector; their correctness is
specified but unverified by a live run. \textit{[AUTHOR THESIS: verification requires Docker
availability and a live OTLP receiver; this is the primary gap in the PARTIAL status.]}

\subsection{Data and Ontology Surfaces}
\label{sec:otel-weaver-exp-005-collector-to-intake-implementation-data}

EXP-005 operates on two ontological surfaces:

\textbf{OTLP span attribute namespace:} The \texttt{process.pi.*} attribute namespace is the
process-intelligence layer within the OTel attribute hierarchy. Key attributes are
\texttt{process.pi.instance\_id} (mandatory, string, namespace gatekeeper) and
\texttt{process.pi.witness.hash} (optional at collector, mandatory at court, 64-char BLAKE3).
These are mapped in \texttt{otel-weaver/mappings/otel-to-pi-witness-map.yaml}.

\textbf{OCEL~2.0 JSON output schema:} The \texttt{file/pi\_json} exporter writes events in
OCEL~2.0-compatible JSON structure. Each event record must carry an \texttt{ocel:type} field
(the lowercased activity name), an \texttt{ocel:timestamp} field (the
\texttt{pi.intake.ingested\_at} ISO~8601 value injected by \texttt{transform/pi}), and an
\texttt{ocel:objects} map binding the event to the process instance identified by
\texttt{process.pi.instance\_id}. The collector-integration-model is documented in
\texttt{otel-weaver/intel/collector-integration-model.yaml}.

\subsection{Templates and Rendered Artifacts}
\label{sec:otel-weaver-exp-005-collector-to-intake-implementation-templates}

The ggen manifest \texttt{live-check-intake.manifest.yaml} is the template-level artifact for
this experiment. It declares the expected output artifacts (including
\texttt{pi\_intake\_events.json} and the BLAKE3 receipt at
\texttt{receipts/live-check-intake-receipt.json}) and their hash expectations. The manifest
is the upstream trigger for the receipt chain: when the live collector run materialises
\texttt{pi\_intake\_events.json}, the ggen system is expected to hash the file, compare
against the manifest declaration, and write the receipt. This receipt-manufacture step has
not been executed. The rendered artifact (\texttt{live-check-intake-receipt.json}) does not
exist. The manifest itself is the only rendered template artifact that exists in the
repository.
